\documentclass[10pt]{article}
\usepackage[T1]{fontenc}
\usepackage[utf8]{inputenc}
\usepackage[french]{babel}
\usepackage{amsmath,amssymb,mathrsfs,mathpazo}
\usepackage{geometry}
\geometry{a4paper,margin=2cm}
\usepackage{graphicx,array,tabularx,tabularray,colortbl,multirow,multicol,booktabs,enumitem,pifont,ulem,fancybox,diagbox,cancel}
\usepackage{pstricks,pstricks-add,pst-plot,pst-tree,pst-node,pst-text,pst-func,pst-math,pst-bspline,pst-3dplot,pst-eucl,pst-solides3d}
\let\pscalculate\undefined
\usepackage{pst-bezier}
\usepackage{tikz,pgfplots,tkz-tab}
\usetikzlibrary{arrows,arrows.meta,patterns,calc,positioning,trees,decorations.markings,shapes.geometric,shapes.misc,shapes.arrows,intersections,decorations.pathmorphing,decorations.pathreplacing,backgrounds,shadows,babel}
\pgfplotsset{compat=1.18}
\usepgfplotslibrary{fillbetween}
\usepackage[np]{numprint}
\usepackage{listings,makecell,pas-tableur,twoopt}
\let\barre\undefined
\usepackage[pstricks]{tablvar}
\definecolor{darkgreen}{RGB}{0,100,0}
\providecommand{\red}{\color{red}}
\providecommand{\blue}{\color{blue}}
\providecommand{\green}{\color{green}}
\providecommand{\black}{\color{black}}
\providecommand{\white}{\color{white}}
\providecommand{\gray}{\color{gray}}
\providecommand{\cyan}{\color{cyan}}
\providecommand{\magenta}{\color{magenta}}
\providecommand{\yellow}{\color{yellow}}
\providecommand{\violet}{\color{violet}}
\providecommand{\orange}{\color{orange}}
\providecommand{\euro}{\texteuro}
\providecommand{\EUR}{\texteuro}
\providecommand{\ieme}{\textsuperscript{e}}
\providecommand{\degres}{\ensuremath{^\circ}}
\providecommand{\degree}{\ensuremath{^\circ}}
\providecommand{\vv}[1]{\overrightarrow{#1}}
\providecommand{\Coord}[1]{#1}
\providecommand{\decosix}{}
\providecommand{\starredbullet}{\ensuremath{\bullet}}
\pagestyle{empty}
\setlength{\parindent}{0pt}
\usepackage{adjustbox}\DeclareUnicodeCharacter{2236}{:}\DeclareUnicodeCharacter{2190}{\ensuremath{\leftarrow}}
\providecommand{\R}{\mathbb{R}}
\providecommand{\N}{\mathbb{N}}
\providecommand{\Z}{\mathbb{Z}}
\providecommand{\Q}{\mathbb{Q}}
\providecommand{\D}{\mathbb{D}}
\providecommand{\C}{\mathbb{C}}
\providecommand{\e}{{\mathrm{e}}}
\providecommand{\vect}[1]{\overrightarrow{#1}}
\providecommand{\vectt}[1]{\overrightarrow{\mathrm{#1}}}
\providecommand{\vvt}[1]{\overrightarrow{\mathrm{#1}}}
\providecommand{\np}[2][]{{#2~#1}}
\providecommand{\vv}[1]{{\overrightarrow{#1}}}
\renewcommand{\d}{{\mathrm{d}}}
\providecommand{\degres}{{^\circ}}
\providecommand{\degre}{{^\circ}}
\providecommand{\blue}{\color{blue}}
\providecommand{\red}{\color{red}}
\providecommand{\enskip}{\quad}
\providecommand{\llbracket}{\lbrack\!\lbrack}
\providecommand{\rrbracket}{\rbrack\!\rbrack}
\providecommand{\euro}{{\text{€}}}
\providecommand{\no}{{\text{n°}}}
\providecommand{\interval}[2]{{[#1;#2]}}
\providecommand{\intervFO}[2]{{[#1;#2[}}
\providecommand{\dotfill}{\ldots}

\providecommand{\ds}{\displaystyle}
\providecommand{\textding}[1]{#1}
\providecommand{\Oij}{\(\mathrm{O};\vec{\imath},\vec{\jmath}\)}
\providecommand{\Oijk}{\(\mathrm{O};\vec{\imath},\vec{\jmath},\vec{k}\)}
\providecommand{\Ouv}{\(\mathrm{O};\vec{u},\vec{v}\)}
\providecommand{\intervFF}[2]{\ensuremath{[#1;#2]}}
\renewcommand{\euro}{\ifmmode\text{\texteuro}\else\texteuro\fi}
\DeclareUnicodeCharacter{00B0}{\ifmmode^{\circ}\else\textdegree\fi}
\DeclareUnicodeCharacter{20AC}{\ifmmode\text{\texteuro}\else\texteuro\fi}
% Énoncés : Éducation nationale. Transcription et corrigés : APMEP, Denis Vergès et contributeurs. Découpage et classement : CoopMaths, Sébastien Lozano, Éric Elter et contributeurs. Adaptation de la présentation : Maths Explorer.
% CC BY-SA : https://coopmaths.fr/www/about/
% https://coopmaths.fr/alea/static/eam/2025/tex/eam_technologique_2025_06_sujet0v1_3_cor.tex
\begin{document}
\noindent\textit{APMEP / CoopMaths. CC BY-SA. Présentation adaptée par Maths Explorer.}\par\medskip
\medskip

\begin{enumerate}
\item 
	\begin{enumerate}
		\item $\bullet~$La courbe $\mathcal{C}$ contient le point de coordonnées (0~;~2), donc $f(2) = 0$.
		
$\bullet~$La tangente $T$ contient les points (0~;~12) et (2~;~0) ; son coefficient directeur égal au nombre dérivé $f'(2)$ est donc $\dfrac{0 - 12}{2 - 0} = \dfrac{12}{2} = - 6$.
		\item L'ordonnée à l'origine est égale à 12 et le coefficient directeur est égal à $- 6$, donc :

$M(x~;~y) \in T \quad \text{si}\quad y = -6x + 12$.
		\item~
		
\begin{center}
\psset{unit=1cm,arrowsize=2pt 3}
\begin{pspicture}(8.5,2)
\psframe(8.5,2)\psline(0,1.5)(8.5,1.5)\psline(2.5,0)(2.5,2)
\uput[u](1.25,1.4,){$x$}\uput[u](2.7,1.4,){$-2$}\uput[u](4.5,1.4,){$0$}
\uput[u](6.5,1.4,){$4$}\uput[u](8.4,1.4,){$6$}
\rput(1.25,0.75){Variations de $f$}
\psline{->}(2.7,0.4)(4.3,1.3,) \psline{->}(4.7,1.3)(6.3,0.4) \psline{->}(6.7,0.4,)(8.3,1.3,)
\uput[d](4.5,1.5){8}\uput[u](6.5,0){$-8$}\uput[d](4.5,1.5){8}
\end{pspicture}
\end{center}
	\end{enumerate}
\begin{enumerate}
	\item 
$f$ est une fonction polynôme dérivable pour tout réel $s$ et :

$f'(x) = 3 \times 0,5x^2 - 2 \times 3x = 1,5x^2 - 6x = 1,5x(x - 4)$.
\item 
On établit le tableau de signes de cette fonction dérivée :

\begin{center}
\psset{unit=1cm,arrowsize=2pt 3}
\begin{pspicture}(8.5,2)
\psframe(8.5,2)\psline(0,1.5)(8.5,1.5)\psline(2.5,0)(2.5,2)
\psline(0,1.)(8.5,1.)\psline(0,0.5)(8.5,0.5)
\uput[u](1.25,1.4,){$x$}\uput[u](2.7,1.4,){$-2$}\uput[u](4.5,1.4,){$0$}
\uput[u](6.5,1.4,){$4$}\uput[u](8.4,1.4,){$6$}
\rput(1.25,1.25){$1,5x$}\rput(3.5,1.25){$-$}\rput(5.5,1.25){$+$}\rput(7.5,1.25){$+$}
\rput(1.25,0.75){$x - 4$}\rput(3.5,0.75){$-$}\rput(5.5,0.75){$-$}\rput(7.5,0.75){$+$}
\rput(1.25,0.25){$f'(x)$}\rput(3.5,0.25){$+$}\rput(5.5,0.25){$-$}\rput(7.5,0.25){$+$}
\end{pspicture}
\end{center}
	\end{enumerate}
\item $f(x) \leqslant -6 x + 12$ sur l'intervalle [0~;~2] par : géométriquement la courbe $\mathcal{C}$ est \textbf{au-dessous} de la tangente $T$ sur l'intervalle [0~;~2].
\end{enumerate}

\bigskip

\end{document}
